\subsubsection{Checkpoint动态}
图~\ref{fig:qwen4accuracy}说明，最终分数不足以概括整个学习过程。
4B Base在Step150时，Block3的MATH500 Avg@8高1.85个百分点（逐点区间$[0.70,3.00]$），
AIME24高2.92个百分点（$[0.42,5.83]$），但AIME25低4.17个百分点。
Step100的MATH500差值则为$-1.40$（$[-2.57,-0.25]$）；Step50仅AMC23有利于Block3。
官方指令版也存在排序变化：Step50三项Avg@8有利于Token，Step150的AIME24同样有利于Token。
这些反转表明干预改变的是学习路径，而非简单添加一个固定性能增量。
我们保留全部checkpoint，不把观测最高分解释为已验证的权重选择规则。

\subsubsection{为什么选择三个token？块长实验}
\label{sec:blockchoice}
表~\ref{tab:blocksizes}直接给出保留$k=3$所依据的探索性扫描。
Block3在四项Avg@8上均高于Block5，Block10则发生严重退化。
相对Token，Block3有三项Avg@8提高，但AIME24下降，Pass@8也有升有降。
因此，三个token是依据实验保留的短窗口默认值，而非普遍最优块长。

\begin{table}[htbp]
\centering\small
\caption{历史块长扫描：Qwen3-1.7B-Base学生、4B-GRPO教师，DAPO训练200步、seed21。单元格为Avg@8 / Pass@8（\%）；不是后续同机配对实验。}
\label{tab:blocksizes}
\begin{tabular}{@{}lrrrr@{}}
\toprule
Method ($k$) & MATH500 & AIME24 & AIME25 & AMC23 \\
\midrule
Token (1) & 69.12 / 88.60 & 13.33 / 26.67 & 6.67 / 16.67 & 36.75 / 62.65 \\
\textbf{Block3 (3)} & 69.27 / 87.80 & 9.58 / 30.00 & 8.75 / 26.67 & 39.91 / 61.45 \\
Block5 (5) & 67.40 / 87.60 & 9.17 / 20.00 & 7.92 / 23.33 & 34.19 / 61.45 \\
Block10 (10) & 0.00 / 0.00 & 0.00 / 0.00 & 0.00 / 0.00 & 0.00 / 0.00 \\
\bottomrule\end{tabular}

\par\smallskip\footnotesize
Token/Block3/Block5使用A800，Block10使用A100；旧评测未显式固定请求seed。
Block5降低rollout显存预算后重启，其分数来自保留的curated摘要；协议详见附录\ref{history:qwen17_sweep}。
\end{table}

更大的块增加交叉信用项并改变裁剪（式\ref{eq:cross}）。
后续Block10复跑严重退化时仍记录有限梯度，最大范数为867.50和865.21；
其内置AMC23判分存在限制（附录\ref{app:history}）。这些观察未分离坍塌的单一原因。
Random3/Sliding3比较分区方式，并非新增固定块长或独立组件消融。

\subsubsection{提示与局部反馈适配}
最初的4B尝试在更新前已重复并触及长度上限，排除了block损失作为这些初始缺陷的成因。
修复同时移除ChatML并分离请求seed，不是聊天边界符的单因素消融。
各模型输入适配probe见附录\ref{app:probes}。

修复后的4B训练触及上限比例为Token 263/6,400（4.11\%）、Block3 247/6,400（3.86\%）。
指令版引擎length停止分别为94/6,400（1.47\%）、77/6,400（1.20\%），未生成thinking标签。
评测中的boxed缺失与length停止差异仍依赖benchmark（附录\ref{app:instruct}）；长度控制不等于正确率。

\subsubsection{信用重分配与训练动态}
4B的200个更新中，Block3符号分歧率中位数为15.57\%，按advantage绝对值加权后为6.68\%。
相对重分配幅度中位数为1.148，是可超过一的幅度比值；Token对应指标按定义为零。
这些量描述干预，不判断反馈是否正确。

图~\ref{fig:qwen4training}显示后期低熵与梯度下降，未出现Block10式退化，说明重分配本身不诊断坍塌。
附录\ref{app:positionfigures}用覆盖mask辅助展示熵与重合位置图；变化的on-policy前缀使其不能作为固定输入因果检验。
分离共享信号、裁剪与损失尺度的作用，仍需组件对照及重复实验（附录\ref{sec:limits}）。
